\section{Evaluation framework}
\label{sec:evaluation}

\subsection{Three case studies, three jobs}
\label{sec:testbeds}

The three systems are not a repeated benchmark; each answers a question the
others cannot.

\begin{description}[leftmargin=2em,style=nextline]
  \item[\textnormal{\textbf{Lorenz-63 --- the model-error axis, including the
    formulation that relaxes H2.}}]
    A stochastic Lorenz-63 system augmented with a wind-forcing variable. Small
    enough that weak-constraint 4D-Var, which explicitly frees a per-step model
    error term, is affordable to run. This matters because weak-constraint
    4D-Var is the classical method \emph{allowed} to know about model error, and
    is therefore the honest ceiling against which H2's cost must be measured.
  \item[\textnormal{\textbf{Lorenz-96 --- the main testbed.}}]
    Two-scale Lorenz-96: eight slow variables coupled to fast variables,
    integrated with RK4 over 500-step assimilation windows. Observations cover a
    24-dimensional subspace of a 40-dimensional state. Every window draws its own
    physical parameters, randomised by $\pm 20\%$, so no method can memorise a
    single setting. This is the only system carrying the full scheme matrix, and
    all cross-class comparisons are made here.
  \item[\textnormal{\textbf{Quasi-geostrophic channel --- representation and
    realism.}}]
    A two-layer Phillips-type channel integrated in potential vorticity $q$ with
    a spectral inversion for the streamfunction $\psi$, forced by a moving-storm
    wind-stress curl. Observations are sparse meridional columns. Its role is
    H3a: the same physics admits two state variables, which is the cleanest
    available test that the representation is not a neutral container.
\end{description}

\begin{table}[t]
\centering
\small
\begin{tabular}{@{}lccccc@{}}
\toprule
& \textbf{DA baselines} & \textbf{DirectUNet} & \textbf{CFM} & \textbf{SDA} & \textbf{DirectUNet$+$SDA} \\
\midrule
Lorenz-63 & \checkmark\ (incl.\ weak-4D-Var) & --- & --- & --- & --- \\
Lorenz-96 & \checkmark & \checkmark & \checkmark & \checkmark & \checkmark \\
QG        & \checkmark\ (EnKF/ETKF) & \checkmark & --- & --- & --- \\
\bottomrule
\end{tabular}
\caption{Scheme coverage. The matrix is deliberately uneven: cross-class claims
are made on Lorenz-96, where every class is present, and the other two systems
carry the specific arguments described above. \todo{Fill the empty cells, or
state in the final version that they are out of scope rather than pending.}}
\label{tab:coverage}
\end{table}

\subsection{Schemes}
\label{sec:schemes}

Five classes, mapped onto Table~\ref{tab:zeroing}:
\textbf{(i)} classical DA --- strong- and weak-constraint 4D-Var, EnKF, ETKF;
\textbf{(ii)} \textbf{DirectUNet}, a single-pass supervised regression
$\Cv \mapsto \hat{\xv}$, i.e.\ a pure $\Psimean$ estimator;
\textbf{(iii)} \textbf{CFM}, a conditional flow trained on the interpolant
\eqref{eq:interpolant};
\textbf{(iv)} \textbf{SDA}, score-based assimilation with an unconditional or
conditioned prior and observation guidance;
\textbf{(v)} \textbf{DirectUNet\,$+$\,SDA}, blending, where the deterministic
estimate supplies $\Psimean$ and the sampler supplies the remaining slots.

Unrolled variational solvers are deliberately excluded. They are a fourth
operator family whose analysis requires the gradient machinery this paper does
not introduce, and they are treated separately.

\subsection{The model-error axis}
\label{sec:s0s1}

\begin{description}[leftmargin=2em,style=nextline]
  \item[\Sz{} (perfect model).] The assimilating model's parameters equal those
    generating the truth for that window.
  \item[\So{} (model error).] The assimilating model's parameters carry an
    additional $\pm 10\%$ bias. The truth is drawn from the same distribution in
    both scenarios; only $\thv^{*}$ changes.
\end{description}

\paragraph{A caveat that must be stated, not cashed in.}
The \Sz{}/\So{} axis prices model error \emph{only for methods that run a forward
model at inference}. A learned estimator receiving no parameters and integrating
nothing is definitionally unaffected, and its $\So/\Sz \approx 1.00$ is not a
robustness result. Three classes must therefore be kept apart
(Table~\ref{tab:classes}), and a model-free ratio is never tabled as cross-class
robustness evidence.

\begin{table}[t]
\centering
\small
\begin{tabular}{@{}p{0.24\textwidth}p{0.36\textwidth}p{0.32\textwidth}@{}}
\toprule
\textbf{class} & \textbf{what \So{} changes for it} & \textbf{a ratio $\approx 1.00$ means} \\
\midrule
classical DA & integrates a biased forward model; the bias compounds through the
integration & n/a --- these degrade by $1.68$--$1.80\times$ \\
model-free learned (DirectUNet, CFM, unconditional SDA) & nothing: no parameters,
no forward model & \emph{definitional}; carries no information \\
conditioned learned (SDA2/SDA3; QG Q2--Q4) & the conditioning inputs
$(\zv^{*},\thv^{*})$ are biased or corrupted & a genuine robustness result \\
\bottomrule
\end{tabular}
\caption{What \So{} means depends on whether, and how, a scheme consumes model
information. This governs every cross-class statement in the paper.}
\label{tab:classes}
\end{table}

\subsection{Metrics}
RMSE is pooled over all windows and timesteps,
$\sqrt{\overline{(\cdot)^2}}$, with the same convention for every scheme.
Explained variance is reported per field. For ensemble schemes the spread/skill
ratio is compared against the finite-ensemble reliability target
$\sqrt{(N-1)/(N+1)}$, and proper ensemble CRPS / energy scores are computed at a
single uniform convention, $N=30$.

\needsrun{Rank histograms are the primary posterior-diagnostic figure for a DA
readership and are not implemented anywhere in the codebase. This remains the
largest new-code item.}
